\begin{abstract}

    \begin{center}
        \large{% Единое место для указания названия работы
Реализация и автоматизация механизма репликации таблиц трансляций и данных NUMA user-space приложений в ядре Linux 6.6
} \\
    \large\textit{Гадеев Дмитрий Русланович} \\[1 cm]
    \end{center}

    %Краткое описание задачи и основных результатов, мотивирующее прочитать весь текст.
    NUMA-эффект --- это деградация производительности приложений в NUMA-архитектурах, возникающая из-за неравномерного доступа к памяти. Его проявления зависят от локальности размещения потоков и страниц памяти и выражаются в изменении задержки доступа и пропускной способности памяти. Проблема NUMA-эффекта давно известна, и у нее есть множество решений, большинство из которых основаны на предположении о том, что удаленный доступ дороже локального.

    Тем не менее, существующие решения не универсальны, к тому же, большинство из них ориентированы на небольшие нагрузки, умещающиеся внутри одной NUMA-ноды. Потенциал проблемы для больших cross-NUMA нагрузок с неоднозначным соответствием между потоками и доступной памятью не полностью исчерпан. В ходе выполнения данной работы был разработан и исследован механизм репликации таблиц трансляций и данных NUMA user-space приложений для решения проблемы NUMA-эффекта больших cross-NUMA нагрузок. Удалось добиться прироста производительности на нескольких тестовых приложениях.

    %\vfill

    %\textbf{Abstract} \\[1 cm]

    %Research and development of machine learning methods

\end{abstract}
